\MathBookSourcePage{242}
\subsection{$L^p$ 估计(非正则系数)}
\label{sec:ch08-lp-estimates-irregular-coefficients}
\setcounter{thmcount}{0}
\setcounter{equation}{0}

上一节的估计建立在更早各节的结果上, 后者给出了梯度误差的最大范数估计. 特别是, 必须对双线性型中那些未必总能满足的系数作若干假设. 本节给出更一般的结论. 为简单起见, 限制在对称情形, 即假设
\MBSourceItem{1}
\begin{equation}
\MathBookFormulaID{formula-ch08-irregular-symmetric-form}
\label{eq:ch08-irregular-symmetric-form}
a(u,v):=\int_\Omega\sum_{i,j=1}^{d}a_{ij}(x)
\frac{\partial u}{\partial x_i}(x)\frac{\partial v}{\partial x_j}(x)\,dx,
\end{equation}
但只假设 $a_{ij}$ 为有界可测系数.

\MBSourceItem{2}
\begin{Proposition}
\label{prop:ch08-irregular-coefficient-near-two}
设 $a(\cdot,\cdot)$ 由式~\eqref{eq:ch08-irregular-symmetric-form} 给出, 其中 $a_{ij}$ 是满足式~\eqref{eq:ch08-pointwise-ellipticity} 的有界可测系数. 则存在与 $h$ 无关的常数 $\alpha<\infty$, $h_0>0$ 和 $\varepsilon>0$, 使得对所有 $0<h<h_0$ 及 $u_h\in V_h$,
\[
\MathBookFormulaID{formula-ch08-irregular-coefficient-near-two}
|u_h|_{W_p^1(\Omega)}\leq\alpha\sup_{0\neq v_h\in V_h}
\frac{a(u_h,v_h)}{|v_h|_{W_q^1(\Omega)}},
\]
只要 $|2-p|\leq\varepsilon$, 其中 $q$ 是 $p$ 的对偶指数, $1/p+1/q=1$.
\end{Proposition}

先说明这一结果与有限元投影估计的关系. 令 $P_hu$ 为 $u$ 关于单元上双线性型 $a(u,v)$ 到 $V_h$ 上的投影, 即 $P_hu$ 是满足
\[
\MathBookFormulaID{formula-ch08-irregular-galerkin-projection}
a(u-P_hu,v)=0
\qquad\forall v\in V_h
\]
的唯一 $V_h$ 中元素.

\MBSourceItem{3}
\begin{Corollary}
\label{cor:ch08-irregular-projection-stability}
在命题~\ref{prop:ch08-irregular-coefficient-near-two} 的假设下, 投影 $P_h$ 在 $W_p^1(\Omega)$ 中稳定, 即存在与 $h$ 和 $u$ 无关的正常数 $C$, 使得
\[
\MathBookFormulaID{formula-ch08-irregular-projection-stability}
\|P_hu\|_{W_p^1(\Omega)}\leq C\|u\|_{W_p^1(\Omega)},
\qquad |2-p|\leq\varepsilon.
\]
\end{Corollary}

这直接由命题~\ref{prop:ch08-irregular-coefficient-near-two} 得到, 因为 Poincaré 不等式~\ref{prop:ch05-poincare-inequality} 和式~\eqref{eq:ch08-irregular-symmetric-form} 给出相应上界.

在余下证明中, 对 $d$ 分量向量函数 $F$ 定义
\[
\MathBookFormulaID{formula-ch08-vector-lp-norm}
\|F\|_{L^p(\Omega)^d}:=\left(\int_\Omega|F(x)|^p\,dx\right)^{1/p},
\]
其中 $|F(x)|$ 表示欧氏长度. 还使用与定义 1.3.7 等价的 Sobolev 半范数
\[
\MathBookFormulaID{formula-ch08-sobolev-seminorm-gradient}
|w|_{W_p^1(\Omega)}:=\|\nabla w\|_{L^p(\Omega)^d}.
\]

命题~\ref{prop:ch08-irregular-coefficient-near-two} 的证明基于 Meyers (1963) 的思想. 先考虑更简单的双线性型所定义的投影 $P_h^*:V\to V_h$:
\[
\MathBookFormulaID{formula-ch08-laplacian-projection}
(\nabla(P_h^*u-u),\nabla v_h)=0
\qquad\forall v_h\in V_h.
\]
前几节的结果直接蕴含
\MBSourceItem{4}
\begin{equation}
\MathBookFormulaID{formula-ch08-laplacian-projection-stability}
\label{eq:ch08-laplacian-projection-stability}
|P_h^*u|_{W_p^1(\Omega)}\leq C^*|u|_{W_p^1(\Omega)},
\qquad1<p\leq\infty.
\end{equation}
由 Simader (1972) 或 Meyers (1963),
\MBSourceItem{5}
\begin{equation}
\MathBookFormulaID{formula-ch08-meyers-continuous-inf-sup}
\label{eq:ch08-meyers-continuous-inf-sup}
|w|_{W_p^1(\Omega)}\leq c_p\sup_{0\neq v\in W_q^1(\Omega)}
\frac{\langle\nabla w,\nabla v\rangle}{|v|_{W_q^1(\Omega)}}
\qquad\forall w\in\mathring W_p^1(\Omega),
\end{equation}
其中 $c_p$ 是 $1/p$ 的对数凸函数, 因而连续, 且显然 $c_2=1$. 对任意 $u_h\in V_h$, 利用 $P_h^*$,
\MBSourceItem{6}
\begin{equation}
\MathBookFormulaID{formula-ch08-discrete-laplacian-inf-sup}
\label{eq:ch08-discrete-laplacian-inf-sup}
|u_h|_{W_p^1(\Omega)}
\leq c_pC_q\sup_{0\neq v_h\in V_h}
\frac{\langle\nabla u_h,\nabla v_h\rangle}{|v_h|_{W_q^1(\Omega)}}
\qquad1<p<\infty,
\end{equation}
其中 $C_p:=\|P_h^*\|_{\mathring W_p^1(\Omega)\to\mathring W_p^1(\Omega)}\leq C^*$. 显然 $C_2=1$. 为证明 $C_p$ 关于 $p$ 在 $2$ 附近一致连续, 将 $P_h^*$ 视为从 $L^p(\Omega)^d$ 到自身的算子. 使用到 $L^p$ 直和空间的算子插值(见第~\MBUnresolvedReference{chap:ch14}{14} 章), 可得对 $P>2$ 的插值界; 既可使用实插值方法(配合适当的第二指标), 也可使用复插值方法(Bergh and Löfström 1976). 因而, 对任意 $\delta>0$, 存在 $\varepsilon>0$, 使得
\MBSourceItem{7}
\begin{equation}
\MathBookFormulaID{formula-ch08-near-two-laplacian-inf-sup}
\label{eq:ch08-near-two-laplacian-inf-sup}
|u_h|_{W_p^1(\Omega)}\leq(1+\delta)
\sup_{0\neq v_h\in V_h}\frac{\langle\nabla u_h,\nabla v_h\rangle}{|v_h|_{W_q^1(\Omega)}}
\qquad\forall|2-p|\leq\varepsilon,
\end{equation}
其中 $1/p+1/q=1$, 且 $\delta,\varepsilon$ 与 $h$ 无关.

现在用扰动论证处理一般情形. 令 $M$ 为使
\MBSourceItem{8}
\begin{equation}
\MathBookFormulaID{formula-ch08-coefficient-upper-bound}
\label{eq:ch08-coefficient-upper-bound}
\sum_{i,j=1}^{d}a_{ij}(x)\xi_i\xi_j\leq M|\xi|^2
\quad\text{对 a.e. }x\in\Omega
\end{equation}
成立的常数. 定义双线性型
\[
\MathBookFormulaID{formula-ch08-perturbation-bilinear-form}
B(v,w):=\langle\nabla v,\nabla w\rangle-\frac1M a(v,w).
\]
由式~\eqref{eq:ch08-pointwise-ellipticity} 和式~\eqref{eq:ch08-coefficient-upper-bound}, 相关矩阵特征值位于 $[0,1-C_a/M]$. 因而由习题~\ref{exr:ch08-17} 和 Hölder 不等式,
\MBSourceItem{9}
\begin{equation}
\MathBookFormulaID{formula-ch08-perturbation-bound}
\label{eq:ch08-perturbation-bound}
|B(u,v)|\leq\left(1-\frac{C_a}{M}\right)|u|_{W_p^1(\Omega)}|v|_{W_q^1(\Omega)}.
\end{equation}
利用恒等式 $\langle\nabla u,\nabla v\rangle=B(u,v)+M^{-1}a(u,v)$, 并结合式~\eqref{eq:ch08-near-two-laplacian-inf-sup} 和式~\eqref{eq:ch08-perturbation-bound}, 得到
\[
\MathBookFormulaID{formula-ch08-irregular-absorption}
\left(\frac1{1+\delta}-\left(1-\frac{C_a}{M}\right)\right)|u_h|_{W_p^1(\Omega)}
\leq\frac1M\sup_{0\neq v_h\in V_h}\frac{a(u_h,v_h)}{|v_h|_{W_q^1(\Omega)}}.
\]
取 $\delta=(\tfrac12M-C_a)/(M-C_a)$, 并按式~\eqref{eq:ch08-near-two-laplacian-inf-sup} 选择 $\varepsilon$, 即完成命题~\ref{prop:ch08-irregular-coefficient-near-two} 的证明.
